\chapter{Taille du tick, files et priorité}\label{ch:mx:tick-size-queues-and-priority}

Depuis octobre 2016, le pilote de la SEC sur la taille du tick a coté environ douze cents petites valeurs par incréments de cinq cents au lieu d'un cent~; deux de ses trois groupes de
quatre cents valeurs ont également négocié dans ces incréments. Les fourchettes se sont élargies, la profondeur affichée a augmenté, et la valeur d'être en tête d'une file a changé pour chaque teneur de marché de
ces valeurs. Le tick est la plus petite unité de prix du marché, et il décide plus qu'il n'y paraît~: savoir si les fournisseurs de liquidité se disputent sur le prix ou sur
le temps, combien de temps durent les files, combien de temps attend un ordre et ce que vaut sa place dans la file. Ce chapitre fait passer le même flux d'ordres sur des grilles d'un, deux,
cinq et dix cents dans le simulateur de bourse, évalue un emplacement dans la file et lit le pilote de la manière dont il a été conçu pour être lu.

\section{Actifs à gros tick et actifs à petit tick}

\begin{definition}[Pas de cotation relatif]\label{def:mx:tick-size-queues-and-priority:relative}
Le \emph{pas de cotation relatif}\index{pas de cotation relatif} d'un instrument est son \omterm{def:mx:the-limit-order-book:tick}{pas de cotation} $\delta_{\mathrm{tick}}$ divisé par son prix, généralement exprimé en points de base.
Un tick d'un cent représente 25 points de base pour une action à quatre dollars et un quart de point de base pour une action à quatre cents dollars.
\end{definition}

\begin{definition}[Actifs à gros tick et actifs à petit tick]\label{def:mx:tick-size-queues-and-priority:large}
Un \emph{actif à gros tick}\index{actif à gros tick} est un actif dont la fourchette vaut presque toujours un tick~: le tick est contraignant, et les fournisseurs de liquidité qui voudraient
coter une fourchette plus serrée n'en ont pas la possibilité. Un \emph{actif à petit tick}\index{actif à petit tick} présente une fourchette de plusieurs ticks et un carnet clairsemé~: le tick
contraint à peine les prix.
\end{definition}

Le \omterm{def:mx:tick-size-queues-and-priority:relative}{pas de cotation relatif} ne suffit pas à dire à quelle catégorie appartient un actif~: ce qui compte, c'est le tick comparé à la fourchette que l'actif aurait sans lui, laquelle
dépend de la volatilité et de l'activité. Harris (1994) a prédit à partir de cette contrainte ce que des études ultérieures ont mesuré~: un tick plus grand élargit les fourchettes
et concentre la profondeur sur moins de prix. Un gros tick transforme la concurrence sur les prix en concurrence sur le temps~: quand le prix ne peut plus s'améliorer, le seul moyen
d'être premier est d'arriver le premier, et c'est là que se trouvent la \omterm{def:mx:tick-size-queues-and-priority:queuevalue}{valeur de la position dans la file} et une grande partie de la course à la vitesse (Yao et Ye, 2018).

Quelle est la taille d'un tick pour un actif donné~? Robert et Rosenbaum (2011) ont répondu à l'aide d'un modèle.

\begin{definition}[Modèle de zone d'incertitude]\label{def:mx:tick-size-queues-and-priority:uz}
Dans le \emph{modèle de zone d'incertitude}\index{modèle de zone d'incertitude} de Robert et Rosenbaum, le prix négocié ne change que lorsque le prix efficient
$P^\ast$ s'éloigne suffisamment du dernier prix négocié~: autour de chaque valeur au milieu d'un tick existe une zone de largeur $2\eta\,\delta_{\mathrm{tick}}$, $0<\eta\le1$, et
le prix négocié passe à un nouveau tick lorsque $P^\ast$ franchit le bord lointain de la zone. Le paramètre $\eta$ mesure l'aversion des participants de marché
aux variations d'un tick~; Dayri et Rosenbaum (2015) appellent $2\eta\,\delta_{\mathrm{tick}}$ la fourchette implicite d'un \omterm{def:mx:tick-size-queues-and-priority:large}{actif à gros tick}.
\end{definition}

Un $\eta$ faible fait rebondir les variations d'un tick d'un sens à l'autre~: pour qu'un nouveau prix s'établisse, $P^\ast$ doit parcourir $\eta\,\delta_{\mathrm{tick}}$ au-delà du milieu du tick,
et le prix a alors de bonnes chances de revenir en arrière. Pour $\eta=\tfrac12$, le prix négocié est un mouvement brownien échantillonné. L'estimateur compte, parmi les variations d'un tick
du prix négocié, les \emph{continuations} $N_c$ (même sens que la variation précédente) et les \emph{alternances} $N_a$ (sens opposé)~:
$\hat\eta=N_c/(2N_a)$. Dayri et Rosenbaum appellent «~tick optimal~» un tick pour lequel $\eta=\tfrac12$ tout en conservant une fourchette d'un tick~; en dessous d'un demi, le tick est trop
grand pour l'actif.

\omcode{code/firm/queuevalue/firm_queuevalue.py}{67}{73}{Le paramètre de zone d'incertitude à partir des continuations et des alternances d'un tick.}\label{lst:mx:tick-size-queues-and-priority:eta}

L'expérience garde fixe le flux d'ordres de \texttt{firm.tape} et ne modifie que la grille~: chaque \omterm{def:mx:the-limit-order-book:orders}{ordre à cours limité} est envoyé à \texttt{firm.exchsim} avec son prix
arrondi en s'éloignant du marché, les cours acheteurs vers le bas et les cours vendeurs vers le haut, à un multiple de un, deux, cinq ou dix cents (voir la \cref{fig:mx:tick-size-queues-and-priority:grid}).
Les fournisseurs de liquidité ne réagissent pas à la nouvelle grille~; l'expérience isole ce que la grille fait aux mêmes ordres.

\omcode{code/microstructure/06-tick-size-queues-and-priority/python/mx_ticks.py}{35}{50}{Le même flux d'ordres sur une grille plus grossière~: chaque prix limite arrondi en s'éloignant du marché, les cours acheteurs vers le bas et les cours vendeurs vers le haut.}\label{lst:mx:tick-size-queues-and-priority:coarse}

\begin{omfigure}
\begin{tikzpicture}[font=\footnotesize,x=0.5cm,y=0.32cm]
\begin{scope}
\node[anchor=west] at (-0.5,11.5) {grille d'un cent};
\foreach \p/\q/\s in {0/2/omDef,1/1/omDef,2/2/omDef,3/3/omDef,4/2/omDef,5/2/omThm,6/3/omThm,7/1/omThm,8/2/omThm,9/2/omThm}{
  \fill[\s!45] (\p,0) rectangle (\p+0.8,\q);
}
\draw[->,gray] (-0.8,0) -- (10.3,0) node[right] {prix};
\foreach \p/\l in {0/96,4/100,9/105}{\node[below] at (\p+0.4,0) {\l};}
\end{scope}
\begin{scope}[xshift=7.2cm]
\node[anchor=west] at (-0.5,11.5) {grille de cinq cents};
\fill[omDef!45] (-1,0) rectangle (-0.2,8);
\fill[omDef!45] (4,0) rectangle (4.8,2);
\fill[omThm!45] (9,0) rectangle (9.8,10);
\draw[->,gray] (-1.3,0) -- (10.3,0) node[right] {prix};
\foreach \p/\l in {-1/95,4/100,9/105}{\node[below] at (\p+0.4,0) {\l};}
\draw[<->] (4.4,10.8) -- (9.4,10.8) node[midway,above] {un tick};
\end{scope}
\end{tikzpicture}

\omcaption{Les mêmes ordres au repos sur deux grilles, prix en cents. L'arrondi en s'éloignant du marché ramène les cours acheteurs de 96 à 99 vers 95 et les cours vendeurs de 101
à 105 vers 105~: la fourchette devient un tick de cinq cents, et la file du meilleur cours vendeur regroupe les ordres de cinq anciens niveaux.}\label{fig:mx:tick-size-queues-and-priority:grid}
\end{omfigure}

Sur quatre heures simulées par grille, la fourchette pondérée par le temps vaut 1.07, 2.06, 5.03 et 10.03 cents, et vaut un tick pendant 94.2~\%, 97.1~\%, 99.4~\% et 99.7~\% du temps.
La taille affichée au meilleur cours, moyennée sur les deux côtés, passe de 1\,740 actions à 3\,020, 6\,560 et 10\,780~; le nombre d'ordres en file à chaque meilleur
cours passe de 8.7 à 15.3, 33.6 et 55.5. Les \omterm{def:mx:the-limit-order-book:orders}{ordres au marché} sont identiques et paient la grille~: une demi-fourchette effective de 0.56, 1.05, 2.53 et 5.02 cents. La
valeur simulée est un \omterm{def:mx:tick-size-queues-and-priority:large}{actif à gros tick} dès un cent~: $\hat\eta$ vaut 0.108 (353 continuations contre 1\,632 alternances), puis 0.058, 0.019 et 0.007
sur les grilles plus grossières. La fourchette implicite $2\hat\eta\,\delta_{\mathrm{tick}}$ vaut 0.22, 0.23, 0.19 et 0.14 cents~: elle est presque une propriété de l'actif, alors que
la \omterm{def:mx:decomposing-the-spread:quoted}{fourchette affichée} est une propriété de la grille.

\begin{dated}{2026-09}{Un demi-cent de tick pour les actions américaines}\label{dat:mx:tick-size-queues-and-priority:rule612}
Les modifications de 2024 de la SEC apportées à la Rule 612 de la Regulation NMS (Release 34-101070, publiée le 8 octobre 2024) fixent un incrément minimal de prix de USD 0.005
pour les cotations et les ordres à partir de USD 1.00 sur les actions NMS dont la \omterm{def:mx:decomposing-the-spread:quoted}{fourchette affichée} moyenne pondérée par le temps est inférieure ou égale à USD 0.015, conservent USD 0.01 pour
les autres, et ne fixent aucun incrément minimal pour les transactions. La mise en conformité, prévue initialement le 3 novembre 2025, a été reportée par les ordonnances d'exemption du 31 octobre 2025 et du 11 juin 2026 au
premier jour ouvré de novembre 2027.
\end{dated}

\section{La valeur de la position dans la file}

\begin{definition}[Valeur de la position dans la file]\label{def:mx:tick-size-queues-and-priority:queuevalue}
La \emph{valeur de la position dans la file}\index{valeur de la position dans la file} d'un \omterm{def:mx:the-limit-order-book:orders}{ordre à cours limité} est son profit espéré par action à partir de l'instant où il entre dans une file~: la probabilité
qu'il soit exécuté, multipliée par la fourchette qu'il capture à l'exécution, moins la sélection adverse qu'il subit ensuite, nets des frais. La valeur de la position dans la file est l'écart
entre les valeurs des ordres en tête et en queue d'une même file.
\end{definition}

Un ordre en tête de file s'exécute en premier et s'exécute contre presque tout \omterm{def:mx:the-limit-order-book:orders}{ordre au marché} qui atteint son prix, qu'il soit informé ou non~; un ordre en queue
ne s'exécute que lorsqu'un flux important ou persistant a consommé tout ce qui le précède, ce qui correspond plus souvent à un flux qui continue de faire bouger le prix. La position modifie les deux
facteurs de la valeur, la probabilité et la sélection adverse conditionnelle, et Moallemi et Yuan (2016) ont bâti un modèle de la valeur d'une position dans la file
à partir d'eux. La valeur d'un emplacement dans la file de naissance-mort du chapitre~1 est
\[
V(n)=P_n\,\bigl(c-a(n)\bigr),\qquad P_n=\frac{\mu}{\mu+\nu}\prod_{k=1}^{n}\frac{\mu+k\theta}{\mu+k\theta+\nu},
\]
avec $n$ ordres devant, des \omterm{def:mx:the-limit-order-book:orders}{ordres au marché} au taux $\mu$, des annulations au taux $\theta$ par ordre devant et $\nu$ pour le nôtre, $c$ la demi-fourchette captée
et $a(n)$ la sélection adverse conditionnelle à une exécution.

Le simulateur mesure les deux facteurs. \texttt{mx\_ticks.queue\_study} suit chaque \omterm{def:mx:the-limit-order-book:orders}{ordre à cours limité} qui rejoint ou améliore le meilleur cours~: les ordres et les actions qui le précèdent
à l'entrée, ses exécutions, le délai avant sa première exécution, et, pour chaque exécution, la capture $\pm(m-p)$ ainsi que le mouvement du milieu 30 secondes plus tard. Sur la grille d'un cent,
un ordre qui ouvre un nouveau meilleur niveau s'exécute dans 54~\% des cas et rapporte 0.57 cent par action postée~; derrière 20 ordres ou plus, il s'exécute dans 6~\% des cas et
rapporte 0.001 cent. Sa sélection adverse par action exécutée passe de $-0.71$ cent en tête, un gain (les ordres qui ouvrent un niveau sont placés par les fournisseurs du côté
que favorise le prix efficient), à 0.48 cent derrière 20 ordres. Sur les grilles plus grossières, la tête vaut davantage~: 0.66, 1.26 et 2.05 cents par action postée à
deux, cinq et dix cents, contre 0.005, 0.094 et 0.188 en queue. La \omterm{def:mx:tick-size-queues-and-priority:queuevalue}{valeur de la position dans la file}, tête moins queue, vaut 0.57, 0.66, 1.16 et 1.86 cents~:
plus le tick est grand, plus il y a à gagner à être premier.

\begin{omfigure}
\begin{tikzpicture}
\begin{groupplot}[group style={group size=2 by 1,horizontal sep=1.5cm},width=6.6cm,height=5cm,
  tick label style={font=\footnotesize},label style={font=\small},grid=major,grid style={gray!25},table/col sep=comma,
  xtick={0,1,2,3,4,5,6,7},xticklabels={0,1,2,3,5,8,12,20},xlabel={ordres devant (début de tranche)}]
\nextgroupplot[title={\small probabilité d'exécution},ymin=0,ymax=0.9,
  legend columns=2,legend style={font=\footnotesize,at={(0.5,-0.3)},anchor=north,draw=none,fill=none}]
\addplot[omDef,very thick,mark=*] table[x=k,y=fill1]{figdata/microstructure/06-tick-size-queues-and-priority/queue.csv};
\addplot[omDef,dashed] table[x=k,y=bd1]{figdata/microstructure/06-tick-size-queues-and-priority/queue.csv};
\addplot[omThm,very thick,mark=square*] table[x=k,y=fill5]{figdata/microstructure/06-tick-size-queues-and-priority/queue.csv};
\addplot[omThm,dashed] table[x=k,y=bd5]{figdata/microstructure/06-tick-size-queues-and-priority/queue.csv};
\legend{1 cent,naissance--mort,5 cents,naissance--mort}
\nextgroupplot[title={\small valeur par action affichée (cents)},ymin=0,ymax=2.2,
  legend columns=3,legend style={font=\footnotesize,at={(0.5,-0.3)},anchor=north,draw=none,fill=none}]
\addplot[omDef,very thick,mark=*] table[x=k,y=value1]{figdata/microstructure/06-tick-size-queues-and-priority/queue.csv};
\addplot[omThm,very thick,mark=square*] table[x=k,y=value5]{figdata/microstructure/06-tick-size-queues-and-priority/queue.csv};
\addplot[omProp,very thick,mark=triangle*] table[x=k,y=value10]{figdata/microstructure/06-tick-size-queues-and-priority/queue.csv};
\legend{tick de 1 cent,tick de 5 cents,tick de 10 cents}
\end{groupplot}
\end{tikzpicture}

\omcaption{Position dans la file sur la plateforme simulée, quatre heures par grille. À gauche~: la part des ordres exécutés selon le nombre d'ordres devant à l'entrée, avec
la formule naissance--mort aux taux mesurés sur les mêmes séances (en tirets). À droite~: capture du spread diminuée de la sélection adverse à 30~secondes, par action affichée.
Données~: \texttt{mx\_ticks.tick\_experiment}.}\label{fig:mx:tick-size-queues-and-priority:queue}
\end{omfigure}

\section{Dynamique des files et délai d'exécution}

Les files s'allongent avec le tick, si bien que les ordres attendent plus longtemps~: le délai médian entre l'entrée et la première exécution, parmi les ordres exécutés, est de 11.3~secondes sur la grille à un cent,
grille, 24.1 à 2~cents, 55.7 à 5~cents et 81.4 à 10~cents~; la part des ordres entrant dans la file qui finissent par être exécutés tombe de 15.3~\% à 12.3~\%, 8.0~\% et 5.6~\%.
La file naissance--mort prédit la forme de la probabilité d'exécution mais non son niveau (la \cref{fig:mx:tick-size-queues-and-priority:queue}, à gauche). Au pas d'un cent,
avec $\mu=0.194$ \omterm{def:mx:the-limit-order-book:orders}{ordres au marché} par seconde et par côté et $\theta=0.063$ annulations par ordre et par seconde mesurées sur les mêmes séances, le modèle donne 0.76
en tête, 0.44 derrière trois ordres et 0.13 derrière vingt, contre 0.54, 0.22 et 0.06 observés. La file du modèle est une ligne qui ne fait que se raccourcir~; dans
les faits, un véritable meilleur niveau peut être dépassé par un meilleur, ce qui renvoie tous ceux qui s'y trouvent au second niveau, et les ordres s'annulent plus vite lorsque le prix
évolue à leur détriment (le chapitre~1 a relevé le même écart pour les ordres dépassés). Un modèle de file pour le trading doit intégrer le fait que le prix s'éloigne.

Le délai d'exécution espéré dans le modèle, notre ordre n'étant jamais annulé, vaut $1/\mu+\sum_{k=1}^n 1/(\mu+k\theta)$~: en tête, l'attente dure un intervalle entre deux \omterm{def:mx:the-limit-order-book:orders}{ordres au marché},
chaque ordre devant ajoute le temps qu'il faut pour être exécuté ou annulé. Le tutoriel calcule la distribution complète avec
\texttt{firm.queues.fill\_time\_cdf} (livre~4, chapitre~8) et la compare aux délais observés.

\section{Expériences sur le pas de cotation}

\begin{definition}[Doubles différences]\label{def:mx:tick-size-queues-and-priority:did}
Un estimateur par \emph{doubles différences}\index{doubles différences} (difference-in-differences) de l'effet d'un traitement compare la variation d'un résultat
entre l'avant et l'après pour les unités traitées à celle des unités non traitées (de contrôle) sur les mêmes dates~: $(\bar y_{T,1}-\bar y_{T,0})-(\bar y_{C,1}-\bar
y_{C,0})$. Elle élimine ce qui a changé pour tout le monde, sous l'hypothèse que les résultats des unités traitées et de contrôle auraient évolué en parallèle sans le
traitement~; dans une régression, c'est le coefficient de l'interaction des indicatrices de traitement et de période postérieure.
\end{definition}

Le dispositif pilote a été conçu pour cette estimation. La SEC a enjoint aux bourses et à la FINRA de soumettre un plan en juin 2014 et l'a approuvé le 6 mai 2015~; il a démarré
le 3 octobre 2016 pour deux ans. Trois groupes test d'environ quatre cents petites valeurs chacun ont été constitués~: le premier était coté au pas de 5~cents et négocié à l'ancien pas,
le deuxième était coté et négocié au pas de 5~cents (avec des exemptions pour les transactions au prix milieu, de détail et négociées de gré à gré), le troisième ajoutait une règle trade-at
qui réservait les exécutions aux plateformes affichant le meilleur cours. Les valeurs éligibles restantes formaient le groupe de contrôle.

Le groupe de contrôle compte, car le marché a changé au même moment~: dans la comparaison faite par les services de la SEC entre juin--septembre 2016 et novembre
2016--février 2017, les fourchettes se sont nettement élargies pour chaque groupe, groupe de contrôle compris. Par rapport à celui-ci, la \omterm{def:mx:decomposing-the-spread:quoted}{fourchette affichée} relative moyenne des groupes test
ont augmenté de cinq à sept points de base (la \omterm{def:mx:decomposing-the-spread:quoted}{fourchette affichée} médiane de l'échantillon était de 18 points de base). La profondeur affichée au meilleur cours a augmenté dans tous les groupes
test, mais les services ont averti que la hausse est largement mécanique, car la profondeur de plusieurs niveaux à un cent se regroupe sur un seul prix à 5~cents~; la profondeur
à moins de 5~cents du meilleur cours n'a pas augmenté significativement dans les deux premiers groupes. Albuquerque, Song et Yao (2020) ainsi que Chung, Lee et Rösch (2020)
ont étudié les effets du dispositif pilote sur les prix et sur les coûts de négociation des petits et des grands ordres.

\begin{omfigure}
\begin{tikzpicture}
\begin{axis}[width=8.5cm,height=5.2cm,xtick={0,1},xticklabels={avant,après},xmin=-0.3,xmax=1.3,ymin=0,ymax=11000,
  ytick={0,2000,4000,6000,8000,10000},scaled y ticks=false,yticklabel style={/pgf/number format/.cd,fixed,1000 sep={\,}},
  ylabel={\small actions au meilleur cours},tick label style={font=\footnotesize},grid=major,grid style={gray!25},legend columns=3,
  legend style={font=\footnotesize,at={(0.5,-0.2)},anchor=north,draw=none,fill=none,/tikz/every even column/.append style={column sep=8pt}},
  table/col sep=comma]
\addplot[omThm,very thick,mark=square*] table[x=period,y=treated]{figdata/microstructure/06-tick-size-queues-and-priority/did.csv};
\addplot[omDef,very thick,mark=*] table[x=period,y=control]{figdata/microstructure/06-tick-size-queues-and-priority/did.csv};
\addplot[omThm,dashed,mark=o] coordinates {(0,2402.8) (1,1998.8)};
\legend{traitées,contrôle,traitées au pas d'origine}
\end{axis}
\end{tikzpicture}

\omcaption{Expérience simulée sur le \omterm{def:mx:the-limit-order-book:tick}{pas de cotation}~: profondeur au meilleur cours pour huit actions traitées et huit actions de contrôle avant et après le passage des actions traitées à une
grille à 5~cents, tandis que chaque action subit davantage de volatilité et de flux d'ordres de bruit. La droite en pointillés est le contrefactuel de tendance parallèle~; à la fin, l'écart avec celle-ci,
est l'estimation par \omterm{def:mx:tick-size-queues-and-priority:did}{doubles différences}. Données~: \texttt{mx\_ticks.did\_panel}.}\label{fig:mx:tick-size-queues-and-priority:did}
\end{omfigure}

Le simulateur peut faire tourner un dispositif pilote dont on connaît la vérité. Seize actions simulées, de liquidité aléatoire, sont observées pendant 20~minutes avant et 20~minutes après, en
trois fenêtres de cinq minutes chacune (96 observations). Après, chaque action subit 50~\% de mouvements du prix efficient en plus et 20~\% d'ordres de bruit en plus~; les huit
actions traitées passent aussi à une grille à 5~cents. Pour la profondeur au meilleur cours (la \cref{fig:mx:tick-size-queues-and-priority:did}), la variation avant--après des actions traitées
est de $+7\,560$ actions, celle du groupe de contrôle de $-400$, et l'estimation par \omterm{def:mx:tick-size-queues-and-priority:did}{doubles différences} de $+7\,960$. Son erreur type groupée par action est de 2\,270, contre
1\,340 si les fenêtres sont supposées indépendantes~: les fenêtres d'une même action partagent sa liquidité, et l'ignorer surestime la précision d'un facteur de
1.7 (erreurs types groupées, livre~4, chapitre~16). Pour la \omterm{def:mx:decomposing-the-spread:quoted}{fourchette affichée}, l'estimation est de $+3.97$ cents, contre une variation avant--après de
$+4.01$~: le choc commun n'a guère déplacé les fourchettes ici, et c'est précisément ce qu'un chercheur ne pourrait savoir sans le groupe de contrôle.

\section{Règles de priorité au-delà du temps}

La priorité prix-temps n'est pas la seule règle (livre~1, chapitre~19). En pro rata, un ordre entrant est réparti entre les ordres au repos au meilleur cours
proportionnellement à leur taille~; avec une règle de l'ordre de tête, l'ordre qui a fixé le prix obtient en premier une part~; certaines plateformes classent les ordres affichés avant les \omterm{def:mx:order-types-and-their-uses:hidden}{ordres cachés}, les lots
standard avant les lots irréguliers, ou les ordres propres d'un courtier en premier. Chaque règle modifie ce que vaut une place dans la file, et donc ce sur quoi les fournisseurs de liquidité se disputent.

Sur la grille à 5~cents, avec le même flux, le pro rata pur aplatit la file~: la part des ordres qui reçoivent une exécution est de 0.78 en tête et de 0.58
derrière 20 ordres ou plus, contre 0.70 et 0.05 sous la priorité temporelle, et le délai médian d'attente pour une première exécution derrière 20 ordres passe de 76.7~secondes à
4.3~secondes. Les exécutions sont partielles, si bien que la valeur par action affichée décroît encore avec la file (0.72 cents en tête, 0.093 au fond), mais la valeur de la position,
écart tête moins fond, est divisée par deux, passant de 1.20 cents à 0.63. Le pro rata récompense la taille plutôt que la vitesse~: un fournisseur qui veut une plus grande part de chaque exécution affiche davantage de quantité
que ce qu'il souhaite négocier, ce qui explique que les marchés à terme en pro rata ajoutent des allocations minimales et des parts réservées à l'ordre de tête.

\begin{omfigure}
\begin{tikzpicture}
\begin{groupplot}[group style={group size=2 by 1,horizontal sep=1.5cm},width=6.6cm,height=4.8cm,
  tick label style={font=\footnotesize},label style={font=\small},grid=major,grid style={gray!25},table/col sep=comma,
  xtick={0,1,2,3,4},xticklabels={0,1,3,8,20},xlabel={ordres devant (début de tranche)}]
\nextgroupplot[title={\small probabilité d'exécution},ymin=0,ymax=0.9,
  legend columns=2,legend style={font=\footnotesize,at={(1.1,-0.3)},anchor=north,draw=none,fill=none}]
\addplot[omDef,very thick,mark=*] table[x=k,y=fill_fifo]{figdata/microstructure/06-tick-size-queues-and-priority/priority.csv};
\addplot[omThm,very thick,mark=square*] table[x=k,y=fill_pr]{figdata/microstructure/06-tick-size-queues-and-priority/priority.csv};
\legend{priorité temporelle,pro rata}
\nextgroupplot[title={\small valeur par action affichée (cents)},ymin=0,ymax=1.4]
\addplot[omDef,very thick,mark=*] table[x=k,y=value_fifo]{figdata/microstructure/06-tick-size-queues-and-priority/priority.csv};
\addplot[omThm,very thick,mark=square*] table[x=k,y=value_pr]{figdata/microstructure/06-tick-size-queues-and-priority/priority.csv};
\end{groupplot}
\end{tikzpicture}

\omcaption{Priorité temporelle face à un pro rata pur sur la grille à cinq centimes, même flux d'ordres, deux heures chacune. Le pro rata répartit les exécutions sur toute la file.
Données~: \texttt{mx\_ticks.priority\_compare}.}\label{fig:mx:tick-size-queues-and-priority:priority}
\end{omfigure}

\section{Tutoriel~: cinq centimes au lieu d'un}\label{tut:mx:tick-size-queues-and-priority:tut}

\begin{tutorial}
\textbf{Objectif.} Faire tourner le même marché sur quatre grilles, valoriser une place dans la file de deux manières et estimer un changement de tick comme l'a fait le pilote.
\textbf{État final~:} les \cref{fig:mx:tick-size-queues-and-priority:queue,fig:mx:tick-size-queues-and-priority:did,fig:mx:tick-size-queues-and-priority:priority}
et les résultats des sections 1 à 5.
\begin{tutsteps}
\item \textbf{Grilles.} \texttt{mx\_ticks.session(mult, seconds, seed)} pour \texttt{mult} valant 1, 2, 5, 10~; \texttt{grid\_stats} pour la fourchette, la profondeur et
le temps d'un tick.
\item \textbf{Files.} \texttt{queue\_study(tape, mult)} et \texttt{firm\_queuevalue.empirical(entries, edges, by)}~; \texttt{rates(tape)} pour $\mu$
et $\theta$, puis \texttt{fill\_probability} et \texttt{firm.queues.fill\_time\_cdf} pour le modèle.
\item \textbf{Tick.} \texttt{eta\_hat} et \texttt{implicit\_spread} sur les prix échangés de chaque grille.
\item \textbf{Priorité.} \texttt{priority\_compare()} fait tourner la grille à cinq centimes sous priorité temporelle et pro rata.
\item \textbf{Expérience.} \texttt{did\_panel()} simule les actions traitées et les actions témoins et ajuste l'interaction avec des erreurs types groupées (clustered)
(statsmodels)~; tracer avec \texttt{fig\_ticks.py}.
\end{tutsteps}
\textbf{Pour aller plus loin.} Laisser les fournisseurs de liquidité réagir à la grille (poster moins d'ordres, mais de plus grande taille, lorsque le tick est grand) et voir lesquels des
effets mécaniques subsistent~; ajouter une règle trade-at sur une seconde plateforme avec le simulateur multi-plateformes de \texttt{firm.exchsim}.
\end{tutorial}

\section{Construction~: la valeur d'une place dans la file}\label{bld:mx:tick-size-queues-and-priority:queuevalue}

\begin{build}
\bfield{Rôle} Ce que vaut une place dans la file, utilisée par la logique de placement du chapitre~17, par le choix de plateforme du chapitre~18 et par les teneurs de marché
du livre~11.
\bfield{Interface} \texttt{fill\_probability(n, mu, theta, nu)}, \texttt{expected\_time\_to\_fill(n, mu, theta)}, \texttt{slot\_value(\dots)}~;
\texttt{ENTRY} et \texttt{empirical(entries, edges, by)} pour les ordres observés~; \texttt{eta\_hat(prices)}, \texttt{implicit\_spread(eta, tick)}.
\bfield{Règles} La file naissance-mort du chapitre~1~; valeurs empiriques par action affichée, exécutions partielles ou totales, sélection adverse mesurée à partir du prix milieu
juste avant chaque exécution~; $\hat\eta$ à partir des seules variations d'un tick du prix échangé.
\bfield{Tests d'acceptation} \texttt{code/firm/queuevalue/tests/}~: la probabilité d'exécution face à une simulation directe de la file~; l'attente moyenne calculée à la main~;
les statistiques par tranche sur des ordres construits à la main~; $\hat\eta$ sur des trajectoires de prix construites à la main (un rebond pur donne zéro).
\bfield{Pour aller plus loin} Un modèle de file dans lequel le meilleur prix peut bouger (un nouveau niveau devant, le niveau vidé), comme chez Moallemi et Yuan~; \omterm{def:mx:tick-size-queues-and-priority:queuevalue}{valeurs de la position dans la file} nettes
des frais du chapitre~7.
\end{build}

\begin{omsources}
\item L.~Harris, «~Minimum price variations, discrete bid-ask spreads, and quotation sizes~», \emph{Review of Financial Studies} 7(1), 1994.
\item C.~Y.~Robert et M.~Rosenbaum, «~A new approach for the dynamics of ultra-high-frequency data: the model with uncertainty zones~», \emph{Journal of
Financial Econometrics} 9(2), 2011.
\item K.~Dayri et M.~Rosenbaum, «~Large tick assets: implicit spread and optimal tick size~», \emph{Market Microstructure and Liquidity} 1(1), 2015.
\item C.~C.~Moallemi et K.~Yuan, «~A model for queue position valuation in a limit order book~», document de travail, 2016.
\item C.~Yao et M.~Ye, «~Why trading speed matters: a tale of queue rationing under price controls~», \emph{Review of Financial Studies} 31(6), 2018.
\item E.~Hu, P.~Hughes, J.~Ritter, P.~Vegella et H.~Zhang, \emph{Tick Size Pilot Plan and Market Quality}, SEC Division of Economic and Risk Analysis
livre blanc, 2018.
\item R.~Albuquerque, S.~Song et C.~Yao, «~The price effects of liquidity shocks: a study of the SEC's tick size experiment~», \emph{Journal of Financial
Economics} 138(3), 2020.
\item K.~H.~Chung, A.~J.~Lee et D.~Rösch, «~Tick size, liquidity for small and large orders, and price informativeness: evidence from the Tick Size Pilot
Program~», \emph{Journal of Financial Economics} 136(3), 2020.
\item U.S. Securities and Exchange Commission, \emph{Regulation NMS: Minimum Pricing Increments, Access Fees, and Transparency of Better Priced Orders},
Release 34-101070, 2024.
\end{omsources}

\section{Exercices}

\begin{exercise}[$\star$]\label{exo:mx:tick-size-queues-and-priority:1}
Quel est le \omterm{def:mx:tick-size-queues-and-priority:relative}{pas de cotation relatif} d'un tick d'un cent sur une action à quatre dollars et sur une action à quatre cents dollars~? Laquelle est la plus susceptible d'être un
\omterm{def:mx:tick-size-queues-and-priority:large}{actif à gros tick}, et que faudrait-il savoir d'autre~?
\end{exercise}

\begin{exercise}[$\star$]\label{exo:mx:tick-size-queues-and-priority:2}
Sur une journée, le prix échangé d'une action enregistre 40 continuations d'un tick et 160 alternances d'un tick. Estimer $\eta$ et la fourchette implicite pour un
tick d'un cent. Le tick est-il trop grand au sens de Dayri et Rosenbaum~?
\end{exercise}

\begin{exercise}[$\star$]\label{exo:mx:tick-size-queues-and-priority:3}
Un ordre rejoint une file où deux ordres sont placés devant lui~; les \omterm{def:mx:the-limit-order-book:orders}{ordres au marché} arrivent à 0.2 par seconde, chaque ordre s'annule à un taux de 0.05 par seconde, le nôtre compris. Quelle est sa
probabilité d'exécution, et son temps moyen d'exécution s'il n'est jamais annulé~?
\end{exercise}

\begin{exercise}[$\star\star$]\label{exo:mx:tick-size-queues-and-priority:4}
Sur une grille à cinq centimes, un ordre en tête de file s'exécute avec une probabilité de 0.7 et subit 0.4 cent de sélection adverse par action exécutée~; un ordre en queue de file s'exécute
avec une probabilité de 0.06 et subit 0.9. Les deux captent 2.5 cents. Quelle est la valeur de chaque place par action, et de combien la tête dépasse-t-elle la queue pour 500 actions~?
\end{exercise}

\begin{exercise}[$\star\star$]\label{exo:mx:tick-size-queues-and-priority:5}
Dans l'expérience simulée, la profondeur au meilleur cours des actions traitées passe de 2\,400 à 9\,960 actions, et celle des actions témoins de 2\,750 à 2\,350. Calculer la
variation avant--après et les \omterm{def:mx:tick-size-queues-and-priority:did}{doubles différences}, et dire quelle hypothèse permet de lire les \omterm{def:mx:tick-size-queues-and-priority:did}{doubles différences} comme un effet causal~?
\end{exercise}

\begin{exercise}[$\star\star$]\label{exo:mx:tick-size-queues-and-priority:6}
Pourquoi l'erreur type de l'estimation de la profondeur est-elle plus grande lorsqu'elle est groupée par action (2\,270) que lorsque les 96 fenêtres sont traitées comme indépendantes (1\,340)~?
Lequel faut-il présenter~?
\end{exercise}

\begin{exercise}[$\star\star\star$]\label{exo:mx:tick-size-queues-and-priority:7}
\emph{Programmation.} Exécuter \texttt{queue\_study} sur la grille à cinq centimes sous priorité temporelle, puis sous pro rata. Comparer la probabilité d'exécution et la valeur par
action affichée en tête de file et derrière 20 ordres, et expliquer pourquoi la valeur continue de baisser avec la file sous pro rata.
\end{exercise}

\begin{exercise}[$\star\star\star$]\label{exo:mx:tick-size-queues-and-priority:8}
\emph{Trouver la faille.} «~Le programme pilote sur le \omterm{def:mx:the-limit-order-book:tick}{pas de cotation} a fait croître la profondeur affichée aux meilleurs cours dans chaque groupe test, donc il a rendu les petites actions plus liquides.~»
\end{exercise}

\section{Problème~: cinq centimes au lieu d'un}

\begin{problem}[{problème du week-end --- cinq centimes au lieu d'un}]\label{pb:mx:tick-size-queues-and-priority:1}
On demande à un régulateur d'élargir le tick des petites actions pour aider les teneurs de marché. Simuler le changement sur le marché simulé et le comparer au pilote.

\medskip\noindent\textbf{Partie I --- La grille.}
\begin{enumerate}
\item Définir un \omterm{def:mx:tick-size-queues-and-priority:large}{actif à gros tick} et dire pourquoi le \omterm{def:mx:tick-size-queues-and-priority:relative}{pas de cotation relatif} ne suffit pas à le déterminer.
\item Qu'est-ce que l'expérience maintient fixe, et qu'est-ce qu'elle change~?
\item Donner la fourchette moyenne et la part de temps à un tick sur les grilles à un centime et à cinq centimes.
\item Comment la taille affichée et le nombre d'ordres au meilleur cours évoluent-ils~?
\item Que paient les \omterm{def:mx:the-limit-order-book:orders}{ordres au marché} sur chaque grille, et pourquoi cela est-il mécanique ici~?
\end{enumerate}

\medskip\noindent\textbf{Partie II --- Files d'attente.}
\begin{enumerate}[resume]
\item Définir la \omterm{def:mx:tick-size-queues-and-priority:queuevalue}{valeur de la position dans la file} et écrire la valeur naissance--mort d'un emplacement.
\item Donner le temps médian avant la première exécution sur les quatre grilles.
\item Comparer les probabilités d'exécution naissance--mort aux probabilités observées à un centime, et expliquer l'écart.
\item Donner la valeur de la tête et celle de la fin de la file sur chaque grille.
\item Pourquoi la sélection adverse est-elle plus forte en fin de file~?
\end{enumerate}

\medskip\noindent\textbf{Partie III --- Le tick.}
\begin{enumerate}[resume]
\item Énoncer l'estimateur par zone d'incertitude de $\eta$.
\item Donner $\hat\eta$ et la fourchette implicite sur les quatre grilles.
\item Que dit la fourchette implicite, presque constante, sur l'actif~?
\item Où s'applique le tick d'un demi-centime prévu par les amendements de 2024~?
\end{enumerate}

\medskip\noindent\textbf{Partie IV --- L'expérience et le verdict.}
\begin{enumerate}[resume]
\item Décrire les trois groupes de test du pilote.
\item Pourquoi les services de la SEC ont-ils eu besoin du groupe témoin~?
\item Que constate-t-il pour les \omterm{def:mx:decomposing-the-spread:quoted}{fourchettes affichées} et pour la profondeur~?
\item Donner les \omterm{def:mx:tick-size-queues-and-priority:did}{doubles différences} simulées pour la profondeur et leurs deux erreurs types.
\item Énoncer le \emph{résultat nommé}~: la variation de la fourchette, de la profondeur et de la \omterm{def:mx:tick-size-queues-and-priority:queuevalue}{valeur de la position dans la file} lorsque le tick est multiplié par cinq sur le marché simulé, à mettre en regard des
effets mesurés dans le pilote.
\item En une phrase~: qui gagne à un tick de cinq centimes~?
\end{enumerate}
\end{problem}

\section{Questions d'entretien}

\begin{interviewq}[$\star$ \iqroles{trader}]\label{iq:mx:tick-size-queues-and-priority:1}
Qu'est-ce qu'un \omterm{def:mx:tick-size-queues-and-priority:large}{actif à gros tick}, et en quoi le trading d'un tel actif diffère-t-il de celui d'un \omterm{def:mx:tick-size-queues-and-priority:large}{actif à petit tick}~?
\end{interviewq}

\begin{interviewq}[$\star\star$ \iqroles{trader, chercheur}]\label{iq:mx:tick-size-queues-and-priority:2}
Pourquoi la position dans la file a-t-elle de la valeur, et quand en vaut-elle davantage~?
\end{interviewq}

\begin{interviewq}[$\star\star$ \iqroles{chercheur}]\label{iq:mx:tick-size-queues-and-priority:3}
Un ordre attend derrière $n$ autres ordres. Écrivez un modèle de sa probabilité d'exécution et dites ce qu'il laisse de côté.
\end{interviewq}

\begin{interviewq}[$\star\star$ \iqroles{chercheur}]\label{iq:mx:tick-size-queues-and-priority:4}
Comment estimeriez-vous l'effet d'une modification du \omterm{def:mx:the-limit-order-book:tick}{pas de cotation} sur les fourchettes à partir d'un pilote avec des actions traitées et des actions témoins~?
\end{interviewq}

\begin{interviewq}[$\star\star$ \iqroles{trader}]\label{iq:mx:tick-size-queues-and-priority:5}
Vous êtes teneur de marché sur un contrat à terme à répartition au prorata. En quoi votre cotation diffère-t-elle d'un marché d'actions à priorité prix-temps~?
\end{interviewq}

\begin{interviewq}[$\star\star\star$ \iqroles{chercheur}]\label{iq:mx:tick-size-queues-and-priority:6}
Comment sauriez-vous si le tick d'une action est trop grand~?
\end{interviewq}
